\chapter{Honest Model Evaluation}

A short report lands on a security analyst's desk. It holds one table: a new intrusion detector,
tested on network traffic captured from a campus network over a month, reaches an accuracy of
$99$ percent. That number is framed in bold type and closed by a single sentence:
``ready to deploy''.

The data contains one attack for every thousand connections. On the same collection, a three-line
program that always answers ``clean'' obtains an accuracy of $99.9$ percent: it reads no feature at
all and never raises a false accusation. The detector trained for hours loses, and its report
contains not one number that leaks that defeat.

What is wrong is the summary: a single number stands in for four numbers with different roles, on
data whose composition is never stated. Chapter 26 explained why test error differs from training
error; this chapter closes three gaps: which metrics are honest about what they ignore, how much
those numbers wobble, and whether the test data is truly unseen by the procedure that chose the
model.

\begin{untukanak}
Imagine a metal detector at a school gate. Every morning it beeps. Yesterday it beeped a thousand
times, and the guard searched a thousand bags. Only sixty of those bags really carried a forbidden
object. The other nine hundred and forty students had to open their bags for no reason.

Now an inspector from the district office arrives and asks: what percentage of the device's answers
were right? The guard answers with pride: ninety-nine percent. He is not lying; out of a hundred
thousand bags, the device was wrong on only a thousand. But that answer does not answer the
question that upsets the students: they are upset not because the device is often wrong, but
because almost every time it beeps, it is wrong.

The two questions sound alike and their answers are far apart: one asks about all the bags, the
other about the bags that were beeped. In this chapter we will learn why one device can excel at
the first question and collapse on the second, why a number of a hundred percent is sometimes not
good news, and why an honest report always carries more than one number.
\end{untukanak}

\section{Four numbers that are not symmetric}
\tingkat{basic}

A binary classifier, however sophisticated its interior, leaves the same trace: four numbers. What
separates one evaluation report from another is whether all four are reported. In everyday
language: every decision has two ways to be right and two ways to be wrong, and closing off one of
them is the same as hiding.

\kamus{confusion matrix}{matriks konfusi}{a four-column ledger that records every decision, right
or wrong, together with the direction of the error}

\begin{definisi}[Confusion matrix]
\index{confusion matrix}
Suppose every example has a true label $y\in\{0,1\}$ and a prediction
$\hat y\in\{0,1\}$, with $1$ meaning positive (attack, disease, harmful message).
The four numbers
\begin{align}
\mathrm{TP}&=\#\{y=1,\ \hat y=1\}, & \mathrm{FP}&=\#\{y=0,\ \hat y=1\},
\nonumber\\
\mathrm{FN}&=\#\{y=1,\ \hat y=0\}, & \mathrm{TN}&=\#\{y=0,\ \hat y=0\}
\nonumber
\end{align}
together with their arrangement are called the \emph{confusion matrix}. The total number of
examples is $n=\mathrm{TP}+\mathrm{FP}+\mathrm{FN}+\mathrm{TN}$.
\end{definisi}

From those four numbers we derive metrics that answer different questions. Precision asks: when the
model says positive, how likely is it to be right. Recall asks: of all the positive cases, what
fraction was caught.

\begin{definisi}[Precision, recall, specificity]
\index{precision}
\index{recall}
With $\mathrm{TP}+\mathrm{FP}>0$ and $\mathrm{TP}+\mathrm{FN}>0$,
\[
\operatorname{prec}=\frac{\mathrm{TP}}{\mathrm{TP}+\mathrm{FP}},\qquad
\operatorname{rec}=\frac{\mathrm{TP}}{\mathrm{TP}+\mathrm{FN}},\qquad
\operatorname{spec}=\frac{\mathrm{TN}}{\mathrm{TN}+\mathrm{FP}} .
\]
Accuracy is $(\mathrm{TP}+\mathrm{TN})/n$, and $F_1$ is
$2\,\mathrm{TP}/(2\,\mathrm{TP}+\mathrm{FP}+\mathrm{FN})$.
\end{definisi}

Let us use one worked example: a detector is tested on $1000$ labeled connections, $100$ of them
attacks. It flags $62$ attacks correctly, misses $38$, and accuses $17$ clean connections of being
attacks.

\begin{center}
\small
\begin{tabular}{l|rr|r}
\hline
 & predicted positive & predicted negative & total \\
\hline
actually positive & $\mathrm{TP}=62$ & $\mathrm{FN}=38$ & 100 \\
actually negative & $\mathrm{FP}=17$ & $\mathrm{TN}=883$ & 900 \\
\hline
total & 79 & 921 & 1000 \\
\hline
\end{tabular}
\end{center}

Its precision is $62/79=0.784810$: about four out of five alarms point to an attack. Its recall is
$62/100=0.620000$: almost four out of ten attacks go unseen. Its specificity is
$883/900=0.981111$. Its accuracy is $945/1000=0.945000$, and $F_1=124/179=0.692737$. All five
numbers are true about the same model, and not one of them may be used on its own. A check with
\texttt{sklearn.metrics} gives the same five numbers to six decimal places.

Figure~\ref{fig:konfusi-jujur} puts the four numbers in their places, together with the row and
column totals; every number in the rest of the chapter must be traceable to one of the cells of
that figure.

\begin{figure}[htbp]
\centering
\begin{tikzpicture}[every node/.style={font=\footnotesize}]
  % two correct cells (green) and two wrong cells (red)
  \fill[layerjade!14] (4.40,2.00) rectangle (6.70,2.90);
  \fill[layerjade!14] (6.70,1.10) rectangle (9.00,2.00);
  \fill[accent!8]     (6.70,2.00) rectangle (9.00,2.90);
  \fill[accent!8]     (4.40,1.10) rectangle (6.70,2.00);
  \draw[ink,thick] (4.40,0.55) rectangle (10.30,2.90);
  \draw[ink] (6.70,0.55) -- (6.70,2.90);
  \draw[ink] (9.00,0.55) -- (9.00,2.90);
  \draw[ink] (4.40,2.00) -- (10.30,2.00);
  \draw[ink] (4.40,1.10) -- (10.30,1.10);
  \draw[ink!35] (4.40,0.55) -- (10.30,0.55);
  % cell contents
  \node at (5.55,2.45) {$\mathrm{TP}=62$};
  \node at (7.85,2.45) {$\mathrm{FN}=38$};
  \node at (5.55,1.55) {$\mathrm{FP}=17$};
  \node at (7.85,1.55) {$\mathrm{TN}=883$};
  % column headers
  \node at (5.55,3.14) {predicted positive};
  \node at (7.85,3.14) {predicted negative};
  \node at (9.65,3.14) {total};
  % row headers
  \node[anchor=east] at (4.24,2.45) {actually positive $(100)$};
  \node[anchor=east] at (4.24,1.55) {actually negative $(900)$};
  \node[anchor=east] at (4.24,0.82) {total};
  % totals
  \node at (5.55,0.82) {$79$};
  \node at (7.85,0.82) {$921$};
  \node at (9.65,2.45) {$100$};
  \node at (9.65,1.55) {$900$};
  \node at (9.65,0.82) {$1000$};
  \node[anchor=north,font=\scriptsize,align=center,text width=11.4cm] at (7.00,0.30)
    {accuracy $=(62+883)/1000=0.945$;\quad precision $=62/79$;\quad recall $=62/100$};
\end{tikzpicture}
\caption{The confusion matrix of the worked example on $1000$ connections: $62$ attacks caught,
$38$ missed, $17$ false alarms, and $883$ clean connections correctly rejected. The rows sum to
$100$ positive and $900$ negative, the columns to $79$ predicted positive and $921$ predicted
negative, and everything to $1000$. The two green cells are the two ways to be right, the two red
cells the two ways to be wrong. Accuracy $0.945$ adds the two green cells and therefore cannot
distinguish the $62$ that were caught from the $38$ that were missed.}
\label{fig:konfusi-jujur}
\end{figure}

The roles of the four numbers are not symmetric. Precision answers the operator's question: when
the alarm sounds, how likely is a real attack. Recall answers the policy designer's question: of
all attacks, what fraction is caught. Both depend on the threshold, so a metric without a
threshold is not a complete statement: the same model gives precision $0.25$ with recall $0.45$,
or precision $0.013$ with recall $0.86$.

$F_1$ tries to compress precision and recall into a single number, and that is exactly where it is
easiest to misuse.

\begin{proposisi}
\index{f1}
$F_1$ is the harmonic mean of precision $P$ and recall $R$, that is $F_1=2PR/(P+R)$,
so that $F_1\le(P+R)/2$ with equality holding exactly when $P=R$.
\end{proposisi}

\begin{proof}[Proof]
Write $a=\mathrm{TP}$, $b=\mathrm{FP}$, $c=\mathrm{FN}$. Then $P=a/(a+b)$,
$R=a/(a+c)$, and
\[
\frac{2PR}{P+R}
=\frac{2a^2\big/\bigl((a+b)(a+c)\bigr)}{a(2a+b+c)\big/\bigl((a+b)(a+c)\bigr)}
=\frac{2a}{2a+b+c}=F_1 .
\]
The harmonic--arithmetic inequality $4PR\le(P+R)^2$ is equivalent to $(P-R)^2\ge0$, and equality
occurs only when $P=R$.
\end{proof}

The harmonic property pulls the value down toward the smaller component: precision $0.9$ with
recall $0.2$ gives $F_1\approx0.327$, not $0.55$; in our example $F_1=0.6927$, slightly below the
midpoint $0.7024$.

Three defects of the summary must be named one by one, because the remedies differ. Precision is
blind to attacks that slip through: a model that flags every connection obtains precision $0.1$,
namely its own prevalence, but recall $1$. $F_1$ ignores $\mathrm{TN}$ entirely: two models with
the same precision and recall but ten times as many false alarms report an identical $F_1$, even
though their operating costs are far apart. $F_1$ also weighs both components as equal, even though
one missed attack is usually far more expensive than one false alarm.

\begin{tantangan}
A report carries only three numbers: ``accuracy $0.9720$ on $5000$ test connections; precision
$0.000$; model ready to deploy''. Determine the whole confusion matrix that is consistent with
that report, explain why its recall is zero and not some other number, then find the largest
prevalence that is still possible.

Hint 1: accuracy only adds the two diagonal cells, so $5000(1-0.9720)=140$ errors are free to be
split between the two remaining cells.

Hint 2: precision $0.000$ means $\mathrm{TP}=0$, so $\mathrm{FN}$ equals the number of positive
examples and $\mathrm{FP}$ is the rest. But precision with $\mathrm{TP}=\mathrm{FP}=0$ is $0/0$,
not $0$; the figure $0.000$ in that report therefore demands at least one false alarm.

If you get stuck, go back to the section Deceptive accuracy on unbalanced data in this chapter.
\end{tantangan}

\section{Deceptive accuracy on unbalanced data}
\tingkat{basic}

Accuracy is the most often reported and the easiest to deceive: it counts all errors at the same
price on a collection whose composition is never stated. In everyday language: accuracy is a report
card computed as if the two groups were balanced.

\index{balanced accuracy}
\begin{definisi}[Balanced accuracy]
Let $\pi$ be the prevalence of the positive class. \emph{Balanced accuracy} is
$(\operatorname{rec}+\operatorname{spec})/2$, that is the average of the correct rates on the two
classes. This quantity does not change when the numbers of positive and negative examples are
changed without changing those two rates.
\end{definisi}

Ordinary accuracy has no such property, and the consequence can be written exactly.

\begin{proposisi}
A classifier that always answers negative obtains accuracy $1-\pi$. Hence a model whose accuracy
is below $1-\pi$ loses to a baseline that learns nothing, and a report that states an accuracy
without stating $\pi$ cannot be judged.
\end{proposisi}

\begin{proof}[Proof]
Of $n$ examples, $\pi n$ are labeled positive and all of them are wrong, while $(1-\pi)n$ are
labeled negative and all of them are right. The accuracy is $(1-\pi)n/n=1-\pi$. Any other model
with smaller accuracy therefore makes more errors than that baseline.
\end{proof}

\begin{bacakeras}
Two voices, read alternately with a friend. Tap the table every time the word ``zero'' is spoken.

``A hundred cards. How many red?'' ``Ten.'' ``How many white?'' ``Ninety.'' ``What does the machine
answer?'' ``White. All white.'' ``How many does it get right?'' ``Ninety.'' ``So its accuracy?''
``Ninety percent.'' ``And the red ones?'' ``Zero. Not a single one.'' ``Was that number lying?''
``No. It answered only one question.'' ``Which question?'' ``How many of everything there is did I
guess right.'' ``Which question went unanswered?'' ``Of the truly dangerous ones, how many did I
find.''

Repeat the three middle lines, faster and faster, with the same taps: what is counted is
everything, what is lost goes unmentioned, what is lost goes unmentioned.

The second voice closes softly: the small class drowns easily inside the average of the large
class, and that average never names who drowned.

The line the reader must repeat, once, in a clear voice: accuracy adds up what is right, not what
matters, so on an unbalanced class ninety percent can mean zero.
\end{bacakeras}

In our worked example $\pi=0.1$, so the baseline is $0.9$ and the model with $0.945$ really is
better. Now change only the prevalence, without changing the quality of the model, and watch what
happens to the same number.

\begin{contoh}[One attack for every thousand connections]
Take $100\,000$ connections with $100$ attacks. Scores on clean connections are normally
distributed with mean $0$ and standard deviation $1$; scores on attacks with mean $2.6$ and
standard deviation $1$. The threshold is set at $2.3263$, the $0.99$ quantile of the negative score
distribution, so by design about one percent of clean connections will be flagged. One run with a
fixed seed gives:
\[
\mathrm{TP}=62,\quad \mathrm{FP}=1013,\quad \mathrm{FN}=38,\quad \mathrm{TN}=98\,887 .
\]
Its accuracy is $(62+98\,887)/100\,000=0.989490$. The baseline is
$99\,900/100\,000=0.999000$. Its precision is $62/1075=0.057674$, its recall
$62/100=0.620000$, and its balanced accuracy $0.804930$.
\end{contoh}

That number $0.9895$, rounded, becomes ``$99$ percent''. It is not merely uninformative: it is
lower than the baseline by $0.009510$, nearly a full percentage point. The model makes ten times
as many errors as the program that always answers ``clean'', $1051$ against $100$, and all of them
are false alarms, $1013$ of them, or about $16.3$ per attack caught.

Here rounding adds damage: $0.9895$ and $0.9990$ both become ``$99$ percent'', even though they
are separated by the whole difference that matters. A good report therefore writes three or four
decimal places, states $n$, and states the prevalence; without all three, the reader does not know
whether $99$ percent is an achievement or a failure.

Should accuracy then not be used at all? It is legitimate when the two errors cost the same, when
the prevalence is near one half, and when the threshold is fixed. Outside that it must be
accompanied by precision and recall, or replaced by balanced accuracy: $0.8049$ is more honest than
$0.9895$, because the minority class is caught only $62$ percent of the time.

\begin{tangan}
Materials: $100$ cards, $90$ white-faced and $10$ red-faced, a marker, and one sheet of tally paper
with two columns.

Write ``clean'' on the ninety white cards and ``attack'' on the ten red ones. Shuffle, then deal
them into ten piles. First rule: every pile is answered ``clean'' without turning over a single
card. Record the correct answers, then turn all the cards over: ninety right, ten red ones through.
Write on the board: accuracy ninety percent, recall zero percent.

Now the second rule. Put a dot in the corner of six red cards, and without telling the player, put
the same dot in the corner of four white cards. Answer ``attack'' only when that dot is visible.
Count again: six red caught, four red through, four white accused, eighty-six white right.

What must become visible: accuracy moves only from ninety to ninety-two percent, while recall jumps
from zero to sixty percent. The number that changes most is precisely the number missing from the
first report.
\end{tangan}

\section{The ROC curve, AUC, and prevalence}
\tingkat{intermediate}

Every threshold gives one point on two different rates; choosing a threshold without looking at all
the possibilities is guessing. In everyday language: we sweep the threshold from loose to strict,
then draw the trail the model leaves behind.

\kamus{receiver operating characteristic}{kurva roc}{the trail a model leaves when its threshold is
swept from loose to strict, drawn on two rates that do not depend on prevalence}

\begin{definisi}[ROC curve and AUC]
\index{roc curve}
\index{auc}
For a threshold $t$, the true positive rate and the false positive rate are
\[
\mathrm{TPR}(t)=P(S>t\mid y=1),\qquad \mathrm{FPR}(t)=P(S>t\mid y=0),
\]
with $S$ the model score. The ROC curve is the set of points $(\mathrm{FPR}(t),\mathrm{TPR}(t))$
for all $t$. The area under that curve is called the AUC, and for finite data it is computed as the
trapezoid area between the points actually observed.
\end{definisi}

The AUC has a far simpler meaning than its name, and that meaning is what made it popular.

\begin{proposisi}
For data without tied scores, the AUC equals the probability that the score of one randomly drawn
positive example exceeds the score of one randomly drawn negative example. When there are tied
scores, each tied pair counts as one half.
\end{proposisi}

\begin{proof}[Proof]
Sort all scores from large to small and sweep the threshold downward. Each time the threshold
passes one positive example, $\mathrm{TPR}$ rises by $1/n_+$ at the height of $\mathrm{FPR}$ at
that moment, namely the fraction of negative examples already passed. Each time the threshold
passes a negative example, $\mathrm{FPR}$ rises at constant $\mathrm{TPR}$ and adds no area. The
total area is therefore
\[
\mathrm{AUC}=\frac1{n_+}\sum_{i:\ y_i=1}\mathrm{FPR}_i
=\frac{1}{n_+n_-}\#\{(i,j):\ y_i=1,\ y_j=0,\ S_i>S_j\},
\]
which is the probability claimed, and that last term is what the Mann--Whitney statistic computes.
Tied scores add no area, so each such pair contributes one half.
\end{proof}

The AUC does not see prevalence: if every negative example is duplicated ten times, the ordering of
scores does not change and the AUC stays. How to compute it on paper is shown in the Worked
examples section.

Now the part that makes the ROC curve dangerous: sweep the threshold in a one-in-a-thousand
simulation.

\begin{center}
\small
\begin{tabular}{rrrrr}
\hline
threshold & recall & FPR & precision & accuracy \\
\hline
3.0000 & 0.450 & 0.00132 & 0.254 & 0.99813 \\
2.8000 & 0.510 & 0.00263 & 0.162 & 0.99688 \\
2.6000 & 0.550 & 0.00467 & 0.105 & 0.99488 \\
2.3263 & 0.620 & 0.01014 & 0.058 & 0.98949 \\
2.0000 & 0.740 & 0.02269 & 0.032 & 0.97707 \\
1.5000 & 0.860 & 0.06729 & 0.013 & 0.93264 \\
\hline
\end{tabular}
\end{center}

On the ROC curve those rows move from $(0.001;0.45)$ to $(0.067;0.86)$: a trail that rises steeply
near the vertical axis, exactly what every practitioner wants to see. Its AUC is $0.966448$. But
the precision column says something else. Even at the strictest threshold, where only $132$ clean
connections are flagged, the precision is $0.254$: three out of four alarms are still wrong. When
recall is raised to $0.86$, the precision is $0.013$.

Figure~\ref{fig:presisi-recall-ambang} draws the six rows of that table on one threshold axis:
recall climbs while precision falls, and no point gives both at once.

\begin{figure}[htbp]
\centering
\begin{tikzpicture}
\begin{axis}[
  width=0.78\linewidth, height=6.2cm,
  x dir=reverse, xmin=1.45, xmax=3.05,
  ymin=0.35, ymax=0.95, ytick={0.4,0.5,0.6,0.7,0.8,0.9},
  xlabel={threshold $t$ (to the right: the threshold is loosened)}, ylabel={recall},
  axis y line*=left, axis x line*=bottom,
  tick label style={font=\footnotesize}, label style={font=\footnotesize},
  grid=major, grid style={layergrey!15},
]
\addplot[accent,thick,mark=*,mark size=1.8pt] coordinates {
  (3.0,0.450) (2.8,0.510) (2.6,0.550) (2.3263,0.620) (2.0,0.740) (1.5,0.860)};
\node[accent,font=\footnotesize,anchor=south west] at (axis cs:2.62,0.665) {recall};
\end{axis}
\begin{axis}[
  width=0.78\linewidth, height=6.2cm,
  x dir=reverse, xmin=1.45, xmax=3.05,
  ymin=0, ymax=0.30, ytick={0,0.1,0.2,0.3},
  axis y line*=right, axis x line*=none, ylabel={precision (right scale)},
  tick label style={font=\footnotesize}, label style={font=\footnotesize},
]
\addplot[layerteal,thick,mark=square*,mark size=1.6pt] coordinates {
  (3.0,0.254) (2.8,0.162) (2.6,0.105) (2.3263,0.058) (2.0,0.032) (1.5,0.013)};
\node[layerteal,font=\footnotesize,anchor=north east] at (axis cs:2.78,0.232) {precision};
\end{axis}
\end{tikzpicture}
\caption{Precision and recall at the six thresholds of the previous table, drawn on one threshold
axis that is turned so that to the right the threshold is loosened. Recall climbs from $0.450$ at
$t=3.0$ to $0.860$ at $t=1.5$, while precision falls from $0.254$ to $0.013$: the two curves move in
opposite directions and there is not one point that gives both at once. The precision scale is on
the right, $0$ to $0.30$; at $t=2.3263$, the threshold used in the simulation, precision $0.058$ with
recall $0.620$.}
\label{fig:presisi-recall-ambang}
\end{figure}

The two readings do not contradict each other: they measure different things.

\begin{proposisi}
The ROC curve and the AUC depend only on the conditional score distributions $S\mid y=1$ and
$S\mid y=0$, so they do not change when the prevalence $\pi$ is changed. Precision at a fixed
threshold does change, and follows
\[
\operatorname{prec}(t)=\frac{\pi\,\mathrm{TPR}(t)}
{\pi\,\mathrm{TPR}(t)+(1-\pi)\,\mathrm{FPR}(t)} .
\]
\end{proposisi}

\begin{proof}[Proof]
By definition, $\mathrm{TPR}$ and $\mathrm{FPR}$ are probabilities conditional on $y$, so neither
contains $\pi$; a curve assembled from the two contains no $\pi$ either. For precision, the law of
total probability gives $P(\hat y=1)=\pi\,\mathrm{TPR}+(1-\pi)\mathrm{FPR}$, while
$P(y=1,\hat y=1)=\pi\,\mathrm{TPR}$. Bayes' rule from Chapter 17 gives the form above.
\end{proof}

The consequence can be measured. If prevalence is changed to one half without changing the
conditional score distributions, the AUC stays at $0.967061$ (the difference $0.0006$ is only
Monte Carlo error), while \emph{average precision} jumps from $0.287636$ to $0.967544$. At a
prevalence of one in a thousand the baseline is $\pi=0.001$, so $0.288$ is still far better than
guessing, but far closer to operational reality than $0.966$.

The rule of thumb is simple: when the positive class is rare, report a precision--recall curve, not
only the ROC curve, and state the prevalence. The ROC curve remains useful for comparing two models
on the same population because it is not polluted by composition; it misleads as soon as it is used
to conclude usefulness on another population or at a particular threshold. The AUC is also blind to
calibration: a model whose scores are twice too large has exactly the same AUC, because the
ordering does not change. Those two defects carry us to the next section.

\section{Calibration: does the number mean a probability}
\tingkat{intermediate}

An operator who receives an alarm needs to know one thing that neither precision nor the AUC
answers: when the model writes $0.8$, are eight of ten such cases really positive. In everyday
language: calibration checks whether the weather forecast that says ``seventy percent rain'' is
really right on seven of ten such days.

\kamus{calibration}{kalibrasi}{the agreement between the number a model promises and what actually
happens in the group sharing that same promise}

\begin{definisi}[Perfect calibration]
\index{calibration}
A score $S\in[0,1]$ is called perfectly calibrated when $P(y=1\mid S=s)=s$ for every possible $s$.
On finite data, calibration is checked by grouping scores into intervals, then comparing the mean
score in that interval with the observed positive rate. The absolute difference, weighted by the
size of the interval, is called the \emph{expected calibration error}.
\end{definisi}

A worked example with $1000$ predictions shows the shape of the check.

\begin{center}
\small
\begin{tabular}{lrrrr}
\hline
score interval & count & positive & mean score & positive rate \\
\hline
$[0.5;0.6)$ & 40 & 21 & 0.55 & 0.525 \\
$[0.6;0.7)$ & 60 & 36 & 0.65 & 0.600 \\
$[0.7;0.8)$ & 120 & 90 & 0.75 & 0.750 \\
$[0.8;0.9)$ & 200 & 176 & 0.85 & 0.880 \\
$[0.9;1.0]$ & 580 & 551 & 0.95 & 0.950 \\
\hline
\end{tabular}
\end{center}

The largest difference is in the second interval, $0.05$, and the overall calibration error is
\[
\tfrac{40}{1000}(0.025)+\tfrac{60}{1000}(0.050)+\tfrac{120}{1000}(0)
+\tfrac{200}{1000}(0.030)+\tfrac{580}{1000}(0)=0.010 .
\]
The direction of the error differs between intervals, but its size is small: one percentage point.
The cross-check is free: the counts in the five intervals add to $1000$, and the positives to $874$.

Calibration is measured by squared error, and there is a neat identity there.

\begin{proposisi}[Brier for calibrated scores]
\index{brier}
Let $S$ be perfectly calibrated and $B=\mathbb{E}[(S-y)^2]$ be the Brier score. Then
\[
B=\mathbb{E}\bigl[S(1-S)\bigr].
\]
In particular the constant score $S=\pi$ is perfectly calibrated with $B=\pi(1-\pi)$, and that
constant score is the least sharp calibrated forecaster.
\end{proposisi}

\begin{proof}[Proof]
Conditional on $S=s$, the variable $y$ takes the value $1$ with probability $s$. Then
\[
\mathbb{E}\bigl[(s-y)^2\mid S=s\bigr]
=s(s-1)^2+(1-s)s^2=s(1-s)\bigl[(1-s)+s\bigr]=s(1-s).
\]
Taking the expectation over $S$ gives the first form. For constant $S=\pi$, the calibration
condition demands $P(y=1)=\pi$, which is indeed true, and the formula gives $\pi(1-\pi)$.
\end{proof}

That identity explains two things: calibration alone is not enough, since the constant score $\pi$
is perfectly calibrated but distinguishes no examples; and for rare events the Brier baseline is
very low, so a model that looks good on accuracy can lose badly on Brier. In our simulation
$\pi(1-\pi)=0.000999$, while the hard labels of the model give a Brier equal to its error rate,
$1051/100\,000=0.010510$. The Brier skill score, $1-B/B_{\mathrm{garis\ dasar}}$, is $-9.52$: the
model is ten times worse than forecasting the constant prevalence.

That negative number is no oddity: hard labels throw away the information about uncertainty, and
for rare events that information is the most valuable of all. A model that outputs $0.06$ is saying
something true, and a report that carries only accuracy throws that sentence away.

\section{Uncertainty of metrics}
\tingkat{advanced}

All the numbers above are computed from a limited number of examples, so all of them wobble when
the data is replaced. In everyday language: a metric without an interval is a single shot with no
indication of how wide the target is.

Chapter 20 already derived the Wilson interval and proved its properties, so we do not repeat it
here. What must be remembered is its form for a proportion $k/n$:
\[
\frac{\hat p+\dfrac{z^2}{2n}\pm
z\sqrt{\dfrac{\hat p(1-\hat p)}{n}+\dfrac{z^2}{4n^2}}}{1+\dfrac{z^2}{n}},
\qquad z=1.96 .
\]

\begin{untukanak}
Two goalkeepers practice stopping penalties. The first stops nine out of ten shots. The second
stops nine hundred out of a thousand shots. The newspaper prints the same headline for both: ninety
percent.

But we feel that the two sentences are not equally strong, and that feeling is right. The first
keeper has faced only ten shots. Perhaps he really is excellent, perhaps he was merely lucky ten
times in a row. The second keeper has faced a thousand shots, and luck over that long a run is
almost impossible. The same number, a different strength.

The question that makes one curious: if we may write only one sentence about the first keeper, may
``ninety percent'' still be written? Yes, provided the sentence is completed. In this chapter we
will learn why the same number means ``between fifty-nine and ninety-eight percent'' on ten shots,
but only ``between eighty-eight and ninety-two percent'' on a thousand shots, and why multiplying
the shots a hundredfold narrows the answer only to a tenth.

\end{untukanak}

The following table applies that formula to accuracy, with $z=1.96$.

\begin{center}
\small
\begin{tabular}{lrrr}
\hline
data & $\hat p$ & Wilson interval & width \\
\hline
$9/10$ & 0.900 & $(0.5959;\ 0.9821)$ & 0.3863 \\
$900/1000$ & 0.900 & $(0.8798;\ 0.9171)$ & 0.0372 \\
$10/10$ & 1.000 & $(0.7225;\ 1.0000)$ & 0.2775 \\
$0/1000$ & 0.000 & $(0.0000;\ 0.0038)$ & 0.0038 \\
\hline
\end{tabular}
\end{center}

The first two rows answer the two goalkeepers. The point estimate is the same, $0.9$, but the first
interval is $0.3863$ wide and the second $0.0372$: ten times narrower, not a hundred times.

The third row is the case that appears most often: not a single error in $10$ trials, and then
``accuracy $100$ percent'' is written. The honest Wilson interval, $(0.7225;1)$, is $0.2775$ wide:
the error rate may still be $27.75$ percent. At the other end, zero events in $1000$ trials gives an
upper bound of $0.003827$, while the rule of three from Chapter 20 gives $0.003$; either is
acceptable, as long as the procedure is named.

Comparing two models by checking whether their intervals overlap is a wrong habit: both are tested
on the same data, and that dependence sharpens the comparison rather than blurring it.

\begin{proposisi}[The difference of two paired models]
Let $b$ be the number of examples that model A gets right and model B gets wrong, and $c$ the
reverse, on the same $n$ examples. Their accuracy difference is $d=(b-c)/n$ with variance
\[
\operatorname{Var}(d)=\frac{1}{n}\Bigl(p_b+p_c-(p_b-p_c)^2\Bigr),
\]
so that its estimated standard error is $\sqrt{b+c-(b-c)^2/n}\,/\,n$.
\end{proposisi}

\begin{proof}[Proof]
Every example falls in one of four cells: right for both, wrong for both, right only for A, right
only for B. Since the examples are independent, $(b,c)$ is multinomial with probabilities $p_b,p_c$.
Then $\operatorname{Var}(b-c)=n\bigl(p_b+p_c-(p_b-p_c)^2\bigr)$, and dividing by $n^2$ gives the
formula above.
\end{proof}

At $b=37$, $c=15$, and $n=1000$, the difference is $0.022$ with standard error $0.007177$, so the
$95$ percent interval is $(0.007932;0.036068)$, entirely above zero, and the exact McNemar test
gives a p-value of $0.003185$. The naive treatment, which assumes the two models independent, gives
the interval $(-0.004296;0.048296)$, which contains zero. The two models can therefore be declared
different, provided the comparison is paired and the test uses the difference, not two separate
intervals.

Finally, the size of the test determines how small a difference can be seen. To detect a difference
$\delta$ at level $\alpha$ and power $1-\beta$,
\[
n\approx\frac{(z_{1-\alpha/2}+z_{1-\beta})^2\,\bigl(p_1(1-p_1)+p_0(1-p_0)\bigr)}
{\delta^2}.
\]
At $p_0=p_1=0.9$, level $0.05$, and power $0.8$, a difference of one point demands
$n\approx14\,128$ examples per model, a difference of two points demands $3532$. A test set of a
thousand examples cannot separate two models one point apart.

\section{Data leakage and nested cross-validation}
\tingkat{advanced}

All the computations above assume one thing that must not be violated: the test number is measured
on data the model has never seen, and never seen either by the procedure that chose the model. In
everyday language: an exam score means something only if the questions were never leaked, including
by the teacher who wrote them.

\kamus{data leakage}{kebocoran data}{exam answers that slip accidentally into the study material, so
that the score goes up without any skill}

\begin{definisi}[Data leakage]
\index{data leakage}
Leakage occurs when information that would not be available at the moment the model is used enters
training, model selection, or reporting, so that the evaluation number measures the ability to
memorize, not the ability to generalize.
\end{definisi}

Its forms are various, and the most frequent ones are invisible. Computing a preprocessing mean or
median from all the data, test data included, is leakage: that statistic should come from training
data only. Selecting the $50$ best features by correlation with the label before separating
training from test is also leakage, since the feature filter has already seen the test labels.
Tuning hyperparameters on the test set and reporting the best number is the leakage most often
committed knowingly. On grouped data, a random split places the same group on both sides; on time
series, it leaks the future into the past. The hardest to detect: a feature that is a consequence
of the label, for example a column ``action taken'' filled in after the incident is known.

Feature-selection leakage can be measured, and its size is startling. Take $200$ randomly labeled
examples and $2000$ features of normal noise, with no signal at all. Select the $20$ features with
the largest absolute correlation with the label, then cross-validate with five folds using a nearest
centroid classifier. If the selection is done once on all the data before the folds are formed, the
accuracy is $0.7800$; if the selection is repeated inside each fold using only its training rows,
$0.4600$. On three other seeds the numbers are $0.775$ and $0.415$, $0.820$ and $0.545$: the leaky
version is always high, the honest one always near $0.5$. That number $0.78$ is an illusion from
beginning to end, and not one line of its code is mistyped.

The protocol is one sentence: separate first, then compute anything. Split the data into training,
validation, and test; compute preprocessing statistics from training only; tune hyperparameters on
validation; touch test once, at the end, for one number. When cross-validation is used, the
preprocessing lives inside the fold. For grouped data, the folds follow the groups; for time data,
they follow the time order.

Figure~\ref{fig:pemisahan-tiga} draws that order as three boxes and one forbidden arrow.

\begin{figure}[htbp]
\centering
\begin{tikzpicture}[every node/.style={font=\footnotesize},
  blok/.style={draw=ink,thick,rounded corners=2pt,align=center,
               minimum width=3.5cm,minimum height=1.5cm,inner sep=3pt}]
  \node[blok,fill=layerteal!12]   (latih) at (0.0,0) {\textbf{train} $60\%$\\ compute statistics};
  \node[blok,fill=layerorange!12] (val)   at (4.8,0) {\textbf{validation} $20\%$\\ choose hyperparameters};
  \node[blok,fill=layerindigo!12] (uji)   at (9.6,0) {\textbf{test} $20\%$\\ touch once};
  \draw[-{Latex},ink,thick] (latih.east) -- (val.west)
    node[midway,above,font=\scriptsize] {retrain};
  \draw[-{Latex},ink,thick] (val.east) -- (uji.west)
    node[midway,above,font=\scriptsize] {locked before inspection};
  \draw[-{Latex},accent,thick,dashed] (uji.south) .. controls (9.6,-1.85) and (0.0,-1.85)
    .. (latih.south);
  \draw[accent,very thick] (4.62,-1.42) -- (4.98,-1.06);
  \draw[accent,very thick] (4.62,-1.06) -- (4.98,-1.42);
  \node[accent,anchor=west,font=\scriptsize] at (5.12,-1.24) {forbidden};
  \node[anchor=north,align=center,text width=11.6cm] at (4.8,-1.98)
    {test statistics, feature filters, and hyperparameter tuning must not cross
     back into the training box};
\end{tikzpicture}
\caption{Three separate data sets, with proportions $60/20/20$; on $1000$ connections that means
$600$, $200$, and $200$. Preprocessing statistics are computed from the training box,
hyperparameters are chosen in the validation box, and the test box is touched once at the end. The
dashed arrow from test back to training is leakage: that single arrow is enough to raise accuracy
from $0.4600$ to $0.7800$ on data that is pure noise.}
\label{fig:pemisahan-tiga}
\end{figure}

What remains is one subtle problem. When cross-validation is used both to choose hyperparameters and
to report quality, the reported number is the best of many trials on the same data, so it is too
optimistic. The way out is two layers of folds.

\kamus{nested cross-validation}{validasi silang bersarang}{two layers of folds, the outer layer to
judge and the inner layer to choose, like two separate examination rooms}

\begin{definisi}[Nested cross-validation]
\index{cross-validation}
\index{nested cross-validation}
Let $\Lambda$ be a set of hyperparameters. The data is split into $K$ outer folds. For each outer
fold $k$: train on the remaining data with every $\lambda\in\Lambda$ using $J$ inner folds, choose
the best $\lambda_k$, retrain on the whole training part of fold $k$ with $\lambda_k$, then measure
on fold $k$. The reported number is the average of those $K$ numbers. Judging and selection
therefore use different data.
\end{definisi}

Its cost is $K$ times higher, and that is why this protocol is often skipped; without the outer
layer, not one number in the report can be read as an estimate of quality on new data.

\section{Why one test is not enough}
\tingkat{intermediate}

The sentence ``we tested the model on $1000$ examples and obtained an accuracy of $0.900$'' sounds
conclusive, yet it carries three uncertainties. In everyday language: one weighing does not tell
you whether the scale wobbles, whether the goods changed, or whether we picked the most pleasant
result.

The first is statistical and already computed: at $n=1000$ the $95$ percent interval is $0.037$
wide, and separating two models one point apart demands $14\,128$ examples. The conclusion ``A is
better'' from a difference that small is a conclusion about a coin toss.

The second is selective, and its name is \emph{winner's curse}. Training $20$ variants of identical
quality and then reporting the best means reporting the maximum of $20$ wobbles: in a simulation
with true accuracy $0.90$ and $n=1000$, the mean of each variant is $0.899987$ while the mean of
their maximum is $0.917337$, so $1.7$ points of that number are luck. When the variants are
dependent, the wobble is smaller, but its direction is always upward.

The third is procedural: a test set used repeatedly stops being a test set. If each comparison has
a $5$ percent chance of deceiving, the chance that at least one of $20$ comparisons deceives is
$1-0.95^{20}=0.6415$, larger than the chance of being right.

The remedies match: report intervals, separate the data used for selection from the data used for
reporting, and lock the test set before the experiment begins. If the test has already been used
repeatedly, the next number is called an exploratory number.

\section{A reporting checklist}
\tingkat{basic}

These seven lines separate a report that can be checked from one that merely convinces. In everyday
language: a receipt does not make the goods arrive, but without a receipt nothing can be traced.

\begin{enumerate}
  \item The size of each class and its prevalence, not only the total $n$.
  \item The four numbers of the confusion matrix, then precision and recall at the threshold
        actually used, together with the value of that threshold.
  \item A confidence interval for every main metric, with the method named; for small $n$ use
        Wilson, not Wald.
  \item A precision--recall curve when the positive class is rare, together with average precision
        and its baseline $\pi$.
  \item Calibration: the score interval bounds, the number of examples per interval, and the
        calibration error.
  \item The separation protocol: what was computed from which data, how many times the test data
        was touched, and how many variants were tried; the best number without that denominator is
        a selected number, not a quality number.
  \item The spread of results across seeds or across folds, together with the code and the random
        seed.
\end{enumerate}

What must not be there: accuracy without prevalence, AUC without a curve, and the phrase ``ready to
deploy''.

\begin{resep}
Four steps for reporting evaluation results that do not mislead:
\begin{enumerate}
  \item Separate the data first: training, validation, test, with their proportions and the
        splitting key, written down before a single computation is run.
  \item Compute preprocessing and select features from training only, tune hyperparameters on
        validation only, then save test for a single touch.
  \item Report all four cells of the confusion matrix together with $n$, $\pi$, and the threshold
        used, then precision, recall, and a confidence interval for every metric.
  \item State how many variants were tried and the spread of results across seeds or across folds,
        so that the reader knows how many degrees of freedom have already been spent.
\end{enumerate}
A warning about order: step 1 must not come after step 2 or step 3. Statistics computed before the
data is separated have already moved test information into training, and after that no number can
repair it; the whole report then measures memory.
\end{resep}

\section{Worked examples}
\tingkat{intermediate}

The three computations below can be repeated on paper.

First, the confusion matrix from the author's own working notes, on $35\,069$ test connections:
$\mathrm{TP}=23\,279$, $\mathrm{FP}=997$, $\mathrm{FN}=590$, $\mathrm{TN}=10\,203$.
\begin{align*}
\operatorname{prec}&=\frac{23\,279}{24\,276}=0.958931, &
\operatorname{rec}&=\frac{23\,279}{23\,869}=0.975282,\\
F_1&=\frac{46\,558}{48\,145}=0.967037, &
\mathrm{acc}&=\frac{33\,482}{35\,069}=0.954746,
\end{align*}
and its balanced accuracy is $0.943132$. One thing in those notes is inconsistent: one sentence
states $11\,200$ attacks and $23\,869$ normal, while the matrix places $23\,869$ in the positive
row. The arithmetic is right, the class labels are swapped, and the prevalence must be read from
the matrix, not from the sentence.

Second, the AUC of a small data set. Positives scoring $0.92$, $0.81$, $0.55$ against negatives
scoring $0.70$, $0.41$, $0.35$, $0.22$ give $12$ pairs; $0.92$ and $0.81$ beat all four, $0.55$
beats three, so $11/12=0.916667$. If one positive score equals one negative score, that pair counts
as one half and the result is $0.875$; \texttt{sklearn} gives the same number when $y$ and the
scores are fed in as vectors.

Third, the effect of prevalence on precision when sensitivity $0.90$ and specificity $0.95$ are held
fixed. With $\operatorname{prec}=\pi(0.9)/\bigl(\pi(0.9)+(1-\pi)(0.05)\bigr)$:

\begin{center}
\small
\begin{tabular}{rrrr}
\hline
$\pi$ & $\operatorname{prec}$ & true alarms & false alarms \\
\hline
0.001 & 0.017699 & 1 & 56 \\
0.010 & 0.153846 & 1 & 5.5 \\
0.100 & 0.666667 & 2 & 1 \\
0.500 & 0.947368 & 18 & 1 \\
\hline
\end{tabular}
\end{center}

The last column is the reciprocal of precision minus one, namely the number of false alarms for
every true alarm. At a prevalence of one in a thousand, one true alarm arrives with $55.5$ false
alarms; at a prevalence of one half, with $0.056$. The same model, the same quality, and a
completely different experience for the operator.

\kertas{Take the matrix of $1000$ predictions from earlier and recompute its five metrics without
looking at the text. Then answer: how many $\mathrm{FP}$ must be removed, without changing
$\mathrm{TP}$ and $\mathrm{FN}$, for precision to become $0.900$? Hint: precision $62/(62+x)=0.9$
gives $x=6.89$, so $\mathrm{FP}$ must fall from $17$ to at most $6$. Check that the accuracy after
that change is $0.956$.}

\section{Computational lab: one number, four ways to read it}
\tingkat{advanced}

Every number in this chapter is computed by one short program, which also shows two things a
formula cannot show: how far leakage raises a number, and how far the winner among many variants is
flattered. The code is run as it stands with \texttt{numpy} and \texttt{scipy}.

\begin{lab}
\begin{verbatim}
import numpy as np
from math import sqrt
from scipy.stats import norm

z = norm.ppf(0.975)


def wilson(k, n):
    p, z2 = k / n, z * z
    mid = (p + z2 / (2 * n)) / (1 + z2 / n)
    half = z * sqrt(p * (1 - p) / n + z2 / (4 * n * n)) / (1 + z2 / n)
    return mid - half, mid + half


def roc_pr(spos, sneg):
    s = np.concatenate([spos, sneg])
    y = np.concatenate([np.ones(len(spos)), np.zeros(len(sneg))])
    o = np.argsort(-s)
    s, y = s[o], y[o]
    tp, fp = np.cumsum(y), np.cumsum(1 - y)
    tpr, fpr = tp / len(spos), fp / len(sneg)
    prec = tp / (tp + fp)
    auc = np.sum(np.diff(np.concatenate([[0], fpr])) *
                 (tpr + np.concatenate([[0], tpr[:-1]])) / 2)
    ap = np.sum(np.diff(np.concatenate([[0], tpr])) * prec)
    return auc, ap


rng = np.random.default_rng(20261025)
Nneg, Npos = 99900, 100
sneg = rng.normal(0.0, 1.0, Nneg)
spos = rng.normal(2.6, 1.0, Npos)
t = 2.3263
tp = int((spos > t).sum()); fp = int((sneg > t).sum())
fn, tn = Npos - tp, Nneg - fp
n = Nneg + Npos
print("TP=%d FP=%d FN=%d TN=%d" % (tp, fp, fn, tn))
print("accuracy %.6f  baseline %.6f" % ((tp + tn) / n, Nneg / n))
print("precision %.6f  recall %.6f" % (tp / (tp + fp), tp / Npos))
print("balanced accuracy %.6f" % ((tp / Npos + tn / Nneg) / 2))
print("brier hard labels %.6f  brier constant %.6f"
      % ((fp + fn) / n, 0.001 * 0.999))
a1, p1 = roc_pr(spos, sneg)
print("pi=0.001: AUC=%.6f  AP=%.6f" % (a1, p1))
a2, p2 = roc_pr(rng.normal(2.6, 1.0, 50000), rng.normal(0.0, 1.0, 50000))
print("pi=0.500: AUC=%.6f  AP=%.6f" % (a2, p2))
for k, m in ((9, 10), (900, 1000), (10, 10), (0, 1000)):
    lo, hi = wilson(k, m)
    print("%4d/%-4d (%.6f, %.6f) width %.6f" % (k, m, lo, hi, hi - lo))
b, c = 37, 15
d = (b - c) / 1000
se = sqrt(b + c - (b - c) ** 2 / 1000) / 1000
print("difference %.6f  se %.6f  CI (%.6f, %.6f)"
      % (d, se, d - z * se, d + z * se))
rt = np.random.default_rng(11)
a = rt.binomial(1000, 0.90, (20000, 20)) / 1000
print("20 variants: mean %.6f  mean max %.6f"
      % (a.mean(), a.max(axis=1).mean()))

rl = np.random.default_rng(4)
nn, dd, kk = 200, 2000, 20
X = rl.normal(size=(nn, dd))
y = rl.permutation(np.repeat([0, 1], nn // 2))
corr = np.array([abs(np.corrcoef(X[:, j], y)[0, 1]) for j in range(dd)])
sel_all = np.argsort(-corr)[:kk]


def cv(pick):
    idx = rl.permutation(nn)
    ok = 0
    for f in np.array_split(idx, 5):
        tr = np.setdiff1d(idx, f)
        cols = pick(tr)
        m1 = X[np.ix_(tr[y[tr] == 1], cols)].mean(0)
        m0 = X[np.ix_(tr[y[tr] == 0], cols)].mean(0)
        d1 = ((X[np.ix_(f, cols)] - m1) ** 2).sum(1)
        d0 = ((X[np.ix_(f, cols)] - m0) ** 2).sum(1)
        ok += (np.where(d1 < d0, 1, 0) == y[f]).sum()
    return ok / nn


def top_in(tr):
    c = np.array([abs(np.corrcoef(X[tr, j], y[tr])[0, 1]) for j in range(dd)])
    return np.argsort(-c)[:kk]


print("feature selection: leaky %.4f  honest %.4f"
      % (cv(lambda tr: sel_all), cv(top_in)))
\end{verbatim}
\end{lab}

The actual output, as it stands:

\begin{verbatim}
TP=62 FP=1013 FN=38 TN=98887
accuracy 0.989490  baseline 0.999000
precision 0.057674  recall 0.620000
balanced accuracy 0.804930
brier hard labels 0.010510  brier constant 0.000999
pi=0.001: AUC=0.966448  AP=0.287636
pi=0.500: AUC=0.967061  AP=0.967544
   9/10   (0.595850, 0.982124) width 0.386274
 900/1000 (0.879848, 0.917091) width 0.037243
  10/10   (0.722467, 1.000000) width 0.277533
   0/1000 (0.000000, 0.003827) width 0.003827
difference 0.022000  se 0.007177  CI (0.007932, 0.036068)
20 variants: mean 0.899987  mean max 0.917330
feature selection: leaky 0.7800  honest 0.4600
\end{verbatim}

Four things must be read from that output. First, the two numbers on the second line must appear
side by side: $0.989490$ and $0.999000$. The direction of the difference decides the whole
interpretation: the model under test is below the baseline.

Second, the AUC barely moves, $0.966448$ against $0.967061$, while AP moves from $0.287636$ to
$0.967544$, a factor of $3.4$: numerical evidence for the previous proposition.

Third, the width of the interval shrinks proportionally to $1/\sqrt n$, from $0.386274$ at $n=10$ to
$0.037243$ at $n=1000$, and $0/1000$ gives an upper bound of $0.003827$.

Fourth, the last two lines carry two diseases that are invisible in a single number: the paired
difference $0.022$ with interval $(0.007932;0.036068)$ does not contain zero, yet
$(0.8798;0.9171)$ and $(0.8563;0.8969)$ overlap widely; and the winner among $20$ identical
variants is reported as $0.917337$ although the true quality is $0.899987$.

The program proves not one proposition in this chapter; it checks claims that have already been
proved. The simulation uses distributions we assume, and honest cross-validation on noise data
gives $0.4600$, not $0.5000$: even a number without signal wobbles.

\begin{untukdosen}
The teaching order I suggest: start from the vendor report in the chapter opening, write $99$
percent and $99.9$ percent on the board, then ask the class to decide which is better. Almost the
whole class picks the detector, and out of that unease the confusion matrix appears as an answer,
not as a definition to be memorized. Once precision and recall exist, then $F_1$ and balanced
accuracy can enter. The ROC curve is best drawn by hand for the six points of the threshold table,
because there the students see for themselves that a point which looks good on the ROC curve can
have precision $0.013$.

Two typical misconceptions. First, the AUC is taken to be a context-free measure of quality, so the
threshold is chosen at random; the symptom shows when students cannot explain the precision at the
threshold in use. Second, cross-validation is taken to be always honest, even though feature
selection or standardization outside the fold is enough to raise the number by $30$ points on data
that is pure noise.

Board activity: split the class into two groups, one computing precision for $\pi=0.001$ and one
for $\pi=0.5$ with the same sensitivity and specificity, then ask both to defend their results. The
first group finds $0.018$ and must explain why the same model feels different. Sign that the class
is ready to go on: they can name one situation in which accuracy $0.99$ is worse than not using a
model at all.
\end{untukdosen}

\section{Summary}
\tingkat{basic}

\begin{ringkasan}
\begin{itemize}
  \item The four numbers of the confusion matrix play unsymmetric roles: precision answers the
        operator's question, recall the policy designer's question, and both depend on the
        threshold; $F_1$ is their harmonic mean and ignores $\mathrm{TN}$.
  \item The baseline accuracy is $1-\pi$: at $\pi=0.001$ a model with $0.989490$ loses to a program
        that always answers negative ($0.999000$), while balanced accuracy $0.804930$ is more honest
        because it does not depend on composition.
  \item The AUC equals the probability that a random positive score exceeds a random negative
        score, so it depends only on the conditional score distributions; it does not change when
        prevalence changes, whereas precision and average precision do, from $0.288$ to $0.968$.
  \item Calibration checks whether a score of $0.8$ means eight out of ten. For calibrated scores,
        Brier equals $\mathbb{E}[S(1-S)]$; the constant score $\pi$ is perfectly calibrated but not
        sharp. Hard labels here give $0.010510$ against $0.000999$.
  \item The width of the interval shrinks proportionally to $1/\sqrt n$: $0.3863$ at $n=10$ and
        $0.0372$ at $n=1000$. To compare two models, use the paired difference: $b=37$, $c=15$
        gives $(0.007932;0.036068)$, whereas the unpaired treatment gives an interval that contains
        zero.
  \item One test is not enough, and leakage raises numbers falsely: $0.7800$ when features are
        selected outside the fold, $0.4600$ when inside; the winner among $20$ identical variants is
        flattered by $1.7$ points; and $20$ comparisons at the $5$ percent level give a chance of
        $0.6415$ that one of them deceives.
\end{itemize}
\end{ringkasan}

\section{Exercises}
\tingkat{intermediate}

\begin{latihan}
\begin{enumerate}
  \item Assess this claim: ``Model A reaches an accuracy of $99.1$ percent and model B $98.7$
        percent on the same test set, so A is better.'' True, false, or not enough evidence? Name
        the part that is wrong and two additional numbers that would strengthen it.
  \item Prove that the Brier score for hard $0/1$ labels equals the error rate,
        $(\mathrm{FP}+\mathrm{FN})/n$. Why does that equality fail for scores that are
        probabilities?
  \item Prove that a constant score $S=s$ minimizes the Brier score exactly when $s=\pi$, with
        minimum value $\pi(1-\pi)$.
  \item Given $\mathrm{TP}=180$, $\mathrm{FP}=45$, $\mathrm{FN}=20$, $\mathrm{TN}=1755$. Compute
        precision, recall, $F_1$, accuracy, specificity, and balanced accuracy to six decimal
        places, then check with \texttt{sklearn.metrics}.
  \item Compute the $95$ percent Wilson interval for $95/100$ and $950/1000$. Compare the two
        widths, then explain why that comparison must not be used to conclude which model is
        better.
  \item Compute the AUC from five positives scoring $0.95$, $0.88$, $0.70$, $0.60$, $0.40$ and five
        negatives scoring $0.75$, $0.65$, $0.50$, $0.30$, $0.10$. Write down the number of correct
        pairs, then check with \texttt{roc\_auc\_score}.
  \item A model has sensitivity $0.80$ and specificity $0.99$. Compute its precision at
        $\pi=0.002$ and $\pi=0.2$, then determine the minimum prevalence for precision to be at
        least $0.5$.
  \item You have $6000$ records from $300$ hosts, about $20$ records each, and $40$ hyperparameter
        settings. Design an evaluation protocol without leakage: state the data split, where
        preprocessing is computed, where hyperparameters are chosen, and how many times test is
        touched.
  \item Assess this claim: ``Because the AUC of this model is $0.95$, the model is useful at any
        threshold we choose.'' True, false, or not enough evidence? Include one supporting
        computation.
\end{enumerate}

Hints for the three hardest items. Item 2: write $(s-y)^2$ for $s\in\{0,1\}$; its value is $1$
exactly on an error. Item 3: differentiate $\mathbb{E}[(s-y)^2]=s^2-2s\pi+\pi$ with respect to
$s$, find its minimum, then use the Brier proposition in this chapter. Item 7: its precision is
$0.8\pi/\bigl(0.8\pi+0.01(1-\pi)\bigr)$; solve
$0.8\pi\ge0.5\bigl(0.8\pi+0.01(1-\pi)\bigr)$ for $\pi$.
\end{latihan}

\begin{untukpeneliti}
What is proved in this chapter: the accuracy of a constant classifier is $1-\pi$; $F_1$ is the
harmonic mean of precision and recall with $F_1\le(P+R)/2$; the AUC equals the probability of a
correct ranking together with its Mann--Whitney form; ROC and AUC are invariant to prevalence while
precision is not; for calibrated scores Brier equals $\mathbb{E}[S(1-S)]$; and the variance formula
for the difference of two paired proportions.

What is popular but not a theorem: that ``AUC $0.9$ means a good model''; that the threshold $0.5$
is neutral; that cross-validation is always honest; and that overlapping intervals mean there is no
difference. The last three are mathematically wrong, not merely incomplete. \emph{Average
precision} is also not an unbiased estimate of the area under the population precision--recall
curve: it is a summary of the observed curve, and its value depends on the sample prevalence.

Primary sources: Wasserman, \emph{All of Statistics}, Chapter 6, Chapters 9--10, Chapter 13 in the
Model Selection part, and Chapter 22 in the part ``Assessing Error Rates and Choosing a Good
Classifier''. Goodfellow, Bengio, and Courville, \emph{Deep Learning}, Chapter 11 in the part
``Performance Metrics'', give the one-in-a-million disease example and the notion of
\emph{coverage}. Athalye, Carlini, and Wagner (2018) show that weak evaluation produces false
security claims: seven of nine defenses without a certificate at ICLR 2018 relied on obfuscated
gradients, and six fell once their evaluation was fixed.

Limits of validity: all simulations use assumed normal distributions, so they measure the behavior
of the procedure, not the truth of the model. The Wilson interval is a score approximation, and
McNemar assumes the discordant pairs are symmetric under the null hypothesis. The leakage result
depends on $n$, $d$, and the number of features selected: $0.78$ is not a constant of nature but
one point on a curve that rises with $d/n$.

Open questions: honest intervals for metrics from mutually dependent folds; calibration tests
without the assumption of uniform score intervals; and reporting uncertainty when the test set is
chosen adaptively.
\end{untukpeneliti}

\peta{Chapter 27 of Wasserman, \emph{All of Statistics}: Chapter 6 (estimators, bias, standard
error), Chapters 9--10 (intervals and tests), Chapter 13 in the Model Selection part, and Chapter 22
in the part ``Assessing Error Rates and Choosing a Good Classifier''. The Wilson interval and the
rule of three are referred from Chapter 20 of this book and from Wilson (1927) and Jovanovic--Levy
(1997). The example of deceptive accuracy follows Goodfellow, Bengio, and Courville, \emph{Deep
Learning}, Chapter 11; the limits of adversarial-defense evaluation follow Athalye, Carlini, and
Wagner (2018), discussed in Chapters 27 and 28. Bayes' rule for precision is proved in Chapter 16;
data splitting in Chapter 24.}
